
\documentclass[12pt,twoside,openright]{book}
\usepackage[utf8]{inputenc}
\usepackage[T1]{fontenc}
\usepackage[spanish,es-noquoting]{babel}
\usepackage[paperwidth=5.5in, paperheight=8.5in, top=24mm, bottom=24mm, inner=24mm, outer=20mm, footskip=12mm]{geometry}
\usepackage{titlesec}
\usepackage{fancyhdr}
\usepackage{setspace}
\usepackage{graphicx}
\graphicspath{ {images/} }
\usepackage{microtype}
\usepackage{csquotes}
\usepackage[scaled=1.10]{ebgaramond}
\usepackage[parfill]{parskip}
\usepackage{xcolor}
\setstretch{1.40}

\begin{document}
\chapter{El Encuentro en Madrid}
Era una tarde de primavera en Madrid. El sol se filtraba entre los edificios de la Gran Vía, creando juegos de luces y sombras que parecían danzar sobre el asfalto. Yo caminaba con la prisa habitual de quien vive en la ciudad, con la mente llena de ruido: correos por responder, facturas, preocupaciones... vibrando bajo, muy bajo.

De repente, en la plaza de Callao, lo vi. No llevaba túnica, ni sandalias, ni tenía un halo brillante sobre la cabeza. Llevaba unos vaqueros desgastados, una camiseta blanca y unas zapatillas cómodas. Estaba sentado en un banco, mirando a la gente pasar con una sonrisa tranquila, una sonrisa que parecía detener el tiempo.

---Hola, Francisco ---dijo cuando pasé cerca.

Me detuve en seco. Nadie me llamaba por mi nombre completo en la calle, y menos un desconocido. Pero al mirarlo a los ojos, supe que no era un desconocido. Había una familiaridad ancestral en su mirada, una paz que desarmaba cualquier defensa.
\end{document}
